Calcul arithmétique. L'assiette vaut (1 $-$ t) puissance e ; les recettes valent donc t (1 $-$ t) puissance e, maximales en t = 1 / (1 + e).\par Lecture : ce sont des courbes de Laffer. Avec une élasticité de 0,5, les recettes sont maximales à un taux de 67 \% ; en dessous, monter le taux rapporte toujours quelque chose, de moins en moins ; au-dessus, il fait baisser les recettes. L'indice 100 marque le maximum de chaque courbe : la figure situe le sommet, elle ne dit rien du montant des recettes.
